\subsection*{Brown's translation and comments: Overview of the appendix}
The complete appendix preserves Brown's earlier English copy as an independent textual state: an introductory note, eleven comments updating the program in 1982, the February 1975 postscript on local systems and integration, and a July extract moving from weak higher structures to stratified topology and reconstruction. The February claims about successive deloopings must be read with the July withdrawal of the simple identification of iterated loop structures with strict commutativity. Brown's comments then distinguish the accomplishments of crossed complexes, group-groupoids, and multiple groupoids from the still broader homotopy program. The July extract explicitly omits the biduality and homology discussion; the corresponding material, added subtitles, and retrospective notes in the 1983 version are not interpolated. The guide's detailed paragraphs identify consequential translation slips and source formulas, while the supplemental Brown concordance accounts for every supplied witness difference. This copy is particularly useful for following the development of the program: its earlier questions coexist with the later author's retractions only when the two versions are compared, not when they are silently merged.

\subsection*{Brown's introductory note (cote 134-2, site 115)}
Brown identifies the following text as an informal translation of the mathematical part of Grothendieck's letter to Breen dated 17--19 July 1975, prepared with assistance from Breen and Loday, and thanks Grothendieck for copies of that letter and the two related earlier letters. This establishes the attribution and the intended relation to the French letters reproduced earlier in cote 134-2; it is not a claim that Brown authored the mathematical program or that all three letters are translated in this one note. The separate translation of the February postscript and the July letter must therefore be distinguished from the scope stated on this cover sheet. The note contains no mathematical correction or change of position. Its lower-left handwritten 116 is archive pagination, not an original manuscript section or a new part of the letter; the archive's site 115 and PDF page 116 accordingly differ by one.

\subsection*{Brown's comment 1: Subsequent work on duality (site 117)}
Brown opens his November 1982 comments by observing that duality has advanced since the 1975 letter and naming work by Begueri and Bester, with the references attributed to Breen. This places the forthcoming translation in its historical sequence: the speculative account of duality in the July letter is not presented as an inventory of what was known in 1982. The pencilled supplements identify a memoir of the Société mathématique de France and a Michigan thesis, but do not supply complete publication details; the darker marginal request to ask Breen for references, with its grouping line, records that incompleteness. The reader therefore should not infer a precise theorem or bibliographical identification from this comment alone. Its role is to qualify the context of the duality program later translated, not to amend any formula in it; no conjecture is withdrawn here.

\subsection*{Brown's comment 2: Gr-categories and symmetric groupoids (site 117)}
The second comment supplies the reference to Sinh's thesis requested by the letter's discussion of group-like categories, adds Saavedra Rivano's work, and places Takeuchi's study of the Villamayor--Zelinsky sequence beside these notions. Brown's literal notation is ``+-groupoid'', accompanied by the description of a coherent symmetric monoidal groupoid with strict associativity and a strict zero object. These strictness qualifications belong to the definition being cited and must be kept when comparing it with the weak structures envisaged by Grothendieck. The handwriting adds the accent to État and replaces an obscured initial with Saavedra before Rivano; the canceled initial remains doubtful. This is bibliographical supplementation of the earlier program, not an assertion that all its higher coherent structures are strict. It prepares the contrast between the strict models in comments 3--5 and the more flexible multiple-groupoid constructions in comment 6.

\subsection*{Brown's comment 3: Group-groupoids and crossed modules (site 117)}
Brown identifies the letter's description of groups in $(1\text{-cat})$ with the equivalence between group-groupoids, explicitly group objects in the category of groupoids, and crossed modules. The Brown--Spencer reference supplies the precise comparison, and the sentence about Verdier reports a belief about prior knowledge in 1965 rather than a documentary priority claim. This is the first concrete algebraic model in the comments and the starting point for the six-category diagram of comment 4. The meaning of group-groupoid here includes the internal group structure; it must not be replaced by an arbitrary groupoid. Nothing in this equivalence alone establishes the full higher homotopy hypothesis. The later warning about strict infinity-groupoids concerns its attempted higher-dimensional extension rather than a retraction of the dimension-two comparison.

\subsection*{Brown's comment 4: Six equivalent algebraic models (sites 117--118)}
The diagram joins poly-T-complexes, cubical T-complexes, simplicial T-complexes, crossed complexes, omega-groupoids, and infinity-groupoids by seven bidirectional equivalences, with letters referring to the following bibliography. It extends the previous comparison and directs the reader to the cube algebra, inductive limit theorems, Ashley's simplicial construction, and Jones's work in preparation. The two vertical connections, the four horizontal connections, and the diagonal from omega-groupoids to crossed complexes are drawn by hand; their object names and reference labels are typed. The omega/infinity distinction is part of Brown's terminology here and is preserved. The heavily overtyped symbol in reference d remains uncertain rather than being supplied from the expected published title. Comment 5 immediately restricts the homotopy types represented by these models, while comment 8 uses their equivalences to construct a nerve. Thus the diagram is a comparison within this strict algebraic setting, not the sought equivalence with all spaces.

\subsection*{Brown's comment 5: The limitation of crossed complexes (site 118)}
Brown explicitly says that the infinity-groupoids in the preceding diagram are inadequate for Grothendieck's purpose. His deliberately approximate explanation describes spaces fibering over a $K(\pi,1)$ with fiber a product of Eilenberg--Mac Lane spaces. The phrase ``roughly speaking'' matters: the comment identifies a restriction in the higher homotopy information carried by the models, without giving a complete necessary-and-sufficient classification theorem with all actions and hypotheses spelled out. This is the key qualification on comment 4 and the reason that merely citing equivalence with crossed complexes does not settle the program of the translated letters. Comment 9 repeats the limitation after the classifying-space construction. The readers preserve the assertion exactly; the guide distinguishes Brown's approximation from a general claim about all spaces or all meanings of infinity-groupoid.

\subsection*{Brown's comment 6: Multiple groupoids and truncated types (sites 118--119)}
Loday's n-cat-groups enter as n-fold categories in the category of groups and are reported to describe truncated homotopy types. This offers a different route from the strict crossed-complex models just restricted in comment 5, and supplies the background for the handwritten van Kampen update in comment 11. The description is brief and does not state a truncation degree or the connectivity conventions. In comparing it with Grothendieck's general n-types, the group-based construction must be read in its connected, pointed setting; the sentence is not by itself a model for arbitrary disconnected spaces. Brown does not retract the crossed-complex equivalences, but changes the proposed algebraic model when greater nonabelian information is required. The literal wording and the n indexing are preserved rather than expanded into a stronger theorem in the readers.

\subsection*{Brown's comment 7: Length-two nonabelian complexes (site 119)}
Brown cites Conduché's work in preparation on the nonabelian chain-complex structure arising from a simplicial group whose normalized chain complex has length two. This directly addresses the letters' search for noncommutative analogues of complexes and provides another concrete development beyond the abelian behavior diagnosed in comment 5. The comment supplies neither the additional operations nor their identities, so it should be used as a pointer rather than as a definition of a general nonabelian derived category. Grothendieck's later note 27 in cote 134-2 records the same report from Brown, giving a direct connection between this bibliographical update and the author's retrospective annotations. Brown's printed title contains ``longeur'', which remains literal in every form; the guide's account does not silently turn that spelling into evidence of a different mathematical notion.

\subsection*{Brown's comment 8: A nerve and the missing categorical extension (site 119)}
Following the equivalences in comment 4, Brown associates a simplicial T-complex $NG$ to an infinity-groupoid $G$ and describes its n-simplices as maps from the homotopy infinity-groupoid of the decorated standard n-simplex into $G$. He relates that object to the corresponding homotopy crossed complex, then states that the analogous nerve for the infinity-categories of the cited reference is still unknown. The four marks below Delta are each a single wave, not the double-wave approximation signs in the initial transcription; their typed-versus-handwritten provenance remains doubtful. Preserving this decoration matters because the formula names a structured simplex rather than asserting an approximate equality. The direction ``to G from'' is plain typed wording, and the later functor arrow points from infinity-categories to simplicial sets. The limitation is a dated statement about this particular extension, not a claim that no nerve of any higher-category model exists. Comment 9 uses the groupoid/crossed-complex nerve that has been constructed.

\subsection*{Brown's comment 9: Classifying spaces and nonabelian cohomology (sites 119--120)}
The two references connect crossed complexes with nonabelian extensions and with homotopy classification of maps. Brown explains that the equivalence with simplicial T-complexes extends the Dold--Kan comparison and yields a classifying space $BC$; the displayed bijection $[X,BC]\cong H^0(X,C)$ is presented beside the Eilenberg--Mac Lane formula $[X,K(A,n)]\cong H^n(X,A)$. The zero on the left-hand theory's cohomological degree is literal and must not be changed to n by analogy: its coefficients are a crossed complex, not merely the abelian group in the preceding formula. The printed hypothesis is that X is a CW-complex. Brown then returns explicitly to comment 5 and calls the results still ``too abelian'', so this construction is an intermediate achievement within the program rather than its completion. His bibliographical ``Astérixes'' and punctuation remain as typed; no later publication data are substituted into the letter.

\subsection*{Brown's comment 10: Algebraic tubular neighborhoods (site 120)}
Brown points to Cox's two papers on algebraic tubular neighborhoods and a paper on homotopy theory of simplicial schemes. These are intended as later developments relevant to the July letter's geometric program, especially the proposed local models and the organization of stratified spaces. The comment itself makes no claim that Cox's constructions satisfy all the precise requirements Grothendieck subsequently sets for dévissage; the author's note 28 earlier in cote 134-2 explicitly reserves judgment because he has not yet seen the work. This distinction prevents a bibliographical connection from being read as a proof of the proposed stratified theory. The punctuation after Math remains visually uncertain. Brown's signature, affiliation, date November 1982, and handwritten U.K. belong to the commentary's provenance and are kept separate from the 1975 letters.

\subsection*{Brown's comment 11: A van Kampen update (site 120)}
The pencilled final item reports that Brown and Loday have proved a van Kampen theorem for 2-cat-groups. It depends on the multiple-groupoid direction introduced in comment 6 and indicates how such algebraic models can support a local-to-global theorem, a central demand of Grothendieck's program. The short note gives no hypotheses, statement, or reference, so the guide does not infer a van Kampen theorem for arbitrary weak higher groupoids from it. Its placement below the date and lack of a relocation cue make it an added numbered comment, preserved as an anchored marginal note, not a sentence to be silently inserted into the 1975 translation. The expression 2-cat-groups is clear on the scan. The added report changes the status of one proposed line of work while leaving the earlier cautions about strict crossed complexes intact.

\subsection*{Brown's February translation: Local systems and reconstruction (sites 122--123)}
The letter dated 17 February 1975 is an afterthought to the preceding homotopy program. It associates a pro-homotopy type with a topos, passes to an ordinary truncated type under local homotopical triviality, and compares cohomological invariance with conditions A and B on locally constant higher stacks. Condition C proposes their description as higher functors from the fundamental n-groupoid, extending the familiar classification of coverings. The qualification ``locally constant'', explicit throughout the surrounding discussion, is missing from C's opening words and must be understood there; arbitrary higher stacks are not classified by the fundamental groupoid alone. Reconstruction from local systems then leads to higher classifying topoi and exact fiber functors. The arrows under lim correctly distinguish arbitrary inductive limits from finite projective limits. The weaker proposal of an ordinary classifying topos is immediately questioned for n=2; the typed marginal comment ``this seems suspect'', with its clipped first letter still doubtful, records the objection. The proposed obstruction through van Kampen and descent over two closed subsets remains a question, not an established counterexample. Brown writes a low 0 after phi in the displayed action on a morphism; its role is composition with the induced fundamental-groupoid map, but the printed subscript is retained. This portion corresponds to Appendix 6 in the earlier French typescript and supplies the reconstruction problem on which the following integration formulas depend; the marginal topoi and Stack glosses do not replace Brown's running terminology.

\subsection*{Brown's February translation: Integration and cohomology (sites 124--126)}
Integration is described as Hom from a final higher stack, and strict Picard structures permit its homotopy groups to express hypercohomology after iterated application of the classifying-object operation. Groupoid-valued stacks supply the initial degree range; a group structure allows one further degree, while the proposed necessity of strict Picard structure for a second delooping is tentative and must be read with the July reconsideration. Several literal features need care. Site 124 calls the object after extension to degree N an n-champ although the construction raises its categorical degree, and prints $N\leq i,n$ where the immediately following sentence requires $N\geq\sup(i,n)$. Its comparison from the classifying object of sections to sections of the classifying stack is called a morphism of $(n-1)$-groupoids, although the continuation asserts isomorphisms through degree $n+1$. The small Gamma subscripts on sites 124 and 125 are crosses, unlike the capital X elsewhere. These are retained and explained rather than repaired as typography. More substantially, Brown's clearly typed ``one knows it fails'' on site 125 does not translate the French ``on a fait ce qu'il fallait'' at site 53, which says that the construction has done what is needed to obtain the canonical comparison; the English phrase reverses the sense of that transition. Conditions A$'$ and B$'$ then restate the proposed characterization using truncated higher functors. B$'$ has the same truncation problem as the French source: for n=0, a map from a point to the one-object groupoid of a nontrivial group induces an equivalence on zero-truncations, but the section map from the fixed points of a free nonempty group-set to that set is not a bijection. The diagram's right lower term also abbreviates the pulled-back coefficient by F. These qualifications concern the proposed mathematics and the translation, not the accuracy of its now-checked literal transcription.

\subsection*{Brown's February translation: Three approaches to cohomology (site 126)}
The closing paragraphs ask whether integration of higher stacks can furnish cohomology without complexes of abelian sheaves, injective resolutions, or semi-simplicial algebra. Hypercoverings suggest a connection with the Čech approach, but independence from semi-simplicial methods is explicitly left open. The resulting comparison is conditional: commutative homological algebra resolves coefficients, homotopical algebra resolves the base, and the higher-stack approach appears to resolve neither. The three descriptions organize the earlier integration proposal and anticipate the main work's later concern with preserving structured homotopy information rather than only derived-category classes. They do not certify a completed third construction. Brown's source distinguishes a handwritten first hacek from a typed v placed over the second C, and keeps the typed closing salutation and Alexandre signature. This portion corresponds to Appendix 8; unlike the consequential omission in the arXiv witness to the French passage, the translated question about avoiding semi-simplicial algebra is present in full.

\subsection*{Brown's July translation: Axioms and the correction about H-spaces (sites 128--129)}
The July extract begins with Grothendieck's response to Breen's constructive proposal. He values its precision but objects to its technical character and, especially, the absence of the geometrically expected functor from simplicial sets to higher groupoids. A sufficiently complete list of desired notions and relations might lead to an explicit construction or an existence-and-uniqueness theorem. Strict Picard n-categories are required to recover complexes of chains of length n, but on site 129 he explicitly withdraws his earlier belief that arbitrarily high loop-space structures correspond simply to commutative topological groups. This correction must accompany the tentative restriction on repeated delooping in the February letter; strict algebraic commutativity and coherent topological commutativity are not interchangeable assumptions. The passage corresponds to Appendix 9 in the revised 1983 version, which adds subtitles and retrospective notes absent from this Brown copy. Brown's visible heading identifies an extract and its July date; two masked title areas cannot be recovered. His body says ``a theorem'', ``functor nerve'', and ``seems to me to be useless'', where the later revised witness uses different wording. These differences mark distinct textual states and do not authorize replacing Brown's translation by the later author's revisions. The next paragraph returns to a nonabelian Dold--Puppe comparison and supplies the concrete group-category model.

\subsection*{Brown's July translation: Group-categories and nonabelian complexes (sites 129--131)}
The nonabelian question asks what replaces the Dold--Puppe description of strict Picard categories when the additive structure is replaced by a group structure. Groups in 1-categories are described explicitly by a homomorphism $d:L_1\to L_0$ and an action of $L_0$ on $L_1$ satisfying equivariance and the Peiffer identity. The resulting cokernel, central kernel, induced action, and Postnikoff--Sinh class in degree three connect this description with crossed modules and with Brown's comments 3 and 7. The higher-dimensional extension is posed as a problem about noncommutative complexes equipped with additional operations; no general formula for those operations is claimed. The intended homotopy interpretation raises the truncation degree by one on passing from topological groups to connected spaces. Brown's copy omits the word ``pointed'' that appears in the revised witness, so the basepoint convention required by this comparison should be supplied in the reader's interpretation, not inserted into the translation. Site 130 has both closing parentheses in $\theta(d(x_1))$; the witness's misplaced parenthesis is a transcription error, while the outer parenthesis's typed-versus-retouched provenance remains uncertain. Brown also writes ``or commutative formal groups'', where the later version has ``and''. Site 131 then explicitly states that a commentary on biduality and homology is omitted. The corresponding discussion in Appendices 10--14 of the 1983 text is therefore absent by design from this earlier copy, and its later subtitle and discussion of the six operations are not interpolated here. The next surviving portion turns to schemes and stratification.

\subsection*{Brown's July translation: The homotopical triangle (sites 132--133)}
The translation distinguishes the rich algebra of a topos from the information retained by its homotopy type, then arranges spaces modulo n-homotopy, n-groupoids, and special n-categories in a triangle of mutual determination. Special, or n-galois, categories are described through locally constant stacks; exact fiber functors are proposed as the means of recovering the fundamental groupoid. This returns to the February reconstruction problem and supplies the locally constant case from which the stratified proposal departs. Brown's site 133 says local n-systems although its definition uses $(n-1)$-stacks; the revised witness's $(n-1)$-systems exposes the index mismatch. The assertion that the étale topos of a scheme is surely not locally homotopically trivial should be read as a warning about general schemes, not a universal exclusion: the étale topos of an algebraically closed field is the elementary point case. Pro-groupoids and ind-categories are proposed to retain the variation over finer coverings when an ordinary homotopy type is insufficient. The three diagram arrows are handwritten, whereas their vertices are typed. The later subtitle is absent, and the stronger or more cautious formulations of the revised witness do not replace Brown's wording. Only the isolated missing word space in examplethe is corrected in the corrected forms, with an explicit apparatus entry.

\subsection*{Brown's July translation: Constructible coefficients and tubular gluing (sites 134--138)}
The next step retains more than the ordinary homotopy type by fixing an equisingular stratification and studying coefficients locally constant on each stratum. Stability under the usual operations motivates the equisingularity condition. In the two-stratum example, sheaves are reconstructed from restrictions and a comparison through the punctured tubular neighborhood; a fibration from the boundary to the closed stratum and an inclusion into the complement provide a geometric description of the gluing functor. Equivalent spans replace the boundary by the punctured tube and then the closed stratum by its tubular neighborhood. This is the content corresponding to Appendix 16, and it supplies the local model for the diagrammatic treatment that follows. Brown's source has several consequential slips. Site 135 calls both complementary strata closed, points the second inclusion from X to $X_1$, prints an extra superscript star on $i_{1*}$, and says that local constancy is true of $X_1$ where the relevant object is the comparison functor applied to $F_1$. Its wording about stability ``not only for a finite partition ... but for all'' obscures the intended distinction between constructibility after refinement and stability for one fixed partition. Site 136 says ``sheaf'' where the revised version says ``scheme'', writes the bundle as $T/X$, and prints a closed interval in the punctured-cone formula, where deleting the collapsed endpoint requires a half-open interval. A neighborhood of $X_0$ in the disjoint open stratum must be interpreted as a punctured tubular object, not an ordinary inclusion of $X_0$ in $X_1$. The diagram's congruence signs are stronger than the homotopy comparisons established in the prose. All these literal features remain in the readers. The isolated is-for-in slip in ``inclusion of the boundary is a manifold'' on site 137 alone receives an apparatus correction. The later Cox note is absent from Brown's copy; his comment 10 supplies the separate historical update.

\subsection*{Brown's July translation: Diagrammatic and fine homotopy types (sites 138--141)}
The tubular span is generalized to a diagram of homotopy types with locally constant coefficients on its vertices and compatibility maps along its arrows. Tame equisingular spaces are proposed to admit a finite dévissage by manifolds with boundary, locally trivial fibrations, and boundary inclusions; adapted stratifications of a map should then express the direct and inverse image operations. Passage to finer stratifications is intended to retain all constructible information in the fine homotopy type. This corresponds to Appendix 17 and anticipates the main text's insistence on homotopy theories of diagrams. The distinction is essential: one needs to localize the category of diagrams, rather than simply replace each object by its image in the ordinary homotopy category and discard coherence. Brown's site 139 obscures that contrast with ``as based'' and the literal ``category of functions'', where the revised version speaks of fractions; the later note explicitly connecting the construction to derivators is absent. Site 138 says homeomorphisms for the compatibility maps of sheaves, although these are homomorphisms and need not be invertible. Site 140 mixes f and its restriction g in the displayed four operations and prints ``non-étale tubular neighbourhoods'', conflicting with the étale construction just proposed. These remain source-level cautions. The finiteness claim on site 141 is part of an explicitly unfinished program: prime-to-characteristic profinite completions of finite polyhedral or compact manifold models are envisaged, with nonliftable positive-characteristic varieties singled out as mysterious. Its stronger present-tense wording should not be read as a proof of the proposed theory. The next paragraph expressly says that a precise algebraic program has not yet been drawn up and turns to linearization.

\subsection*{Brown's July translation: Linearization and reconstruction (sites 141--143)}
The closing mathematical discussion asks how much cohomological information reconstructs a homotopy type: twisted coefficients, exact-sequence functoriality, products, and other cohomology operations lead to the candidate of a triangulated derived category with derived tensor product and Hom. Grothendieck expresses uncertainty about that abelian candidate and greater confidence in reconstruction from the full higher categories of stacks. The stratified version should similarly recover mixed homotopy information from coefficients subordinate to a fixed stratification and the operations along inclusions of unions of strata. This is the precursor to Appendix 18 and to the later demand for a richer notion than a bare triangulated category. Brown's text must be read with its limitations: site 142 says n-stacks of constructible sheaves where the ordinary reconstruction discussion requires a locally constant qualification; the following paragraph separately imposes local constancy on each stratum. The blank before product is empty and is not filled with the witness's cup. Most significantly, the later author's footnote reporting a negative answer to the proposed abelian recovery is absent from this earlier copy, so Brown's question remains open in its own textual state. Site 143 prints operations of fibrations and cofibrations and retains a full NB about Quillen's early-1970s approach; the revised 1983 version changes these to fibers and cofibers, cancels the NB, and adds the derivator discussion. These are distinct conceptual and chronological layers. The letter closes with the request for the Bloch--Quillen formula, not with a claimed solution to the reconstruction program.

